\subsubsection{The Hopf pairing and its balanced section}
\label{sec:Hopf-pairing}

The Hopf pairing realizes weights $(1,8,8)$ and has two linear
singular strata. Its balanced section also gives an intrinsic
comparison inside a Simons cone.

Let \(h:\R^{16}\to\R^9\) be the quadratic extension of the standard
octonionic Hopf map \(\Sph^{15}\to\Sph^8\), normalized by
\begin{equation}\label{eq:Hopf-normalization}
 |h(x)|^2=|x|^4,\qquad
 Dh(x)Dh(x)^{\mathsf T}=4|x|^2\Id_9,\qquad \Delta h=0.
\end{equation}
For polarization of the spherical norm identity and the Hopf construction
from norm-preserving bilinear multiplication, see
\cite[Chapter~15]{Shapiro2000}.
The coordinates are \(h_i(x)=\langle A_ix,x\rangle\) for a nine-matrix Clifford
system on \(\R^{16}\); see \cite[Theorem~A and Section~1.2]{Radeschi14}.
On \(\R^{16}\oplus\R^{16}\), set
\begin{equation}\label{eq:Hopf-pairing-definitions}
 u=|x|^2,\quad v=|y|^2,\quad R=u+v,\quad V=u-v,\quad T=uv,
 \qquad F(x,y)=\langle h(x),h(y)\rangle.
\end{equation}
In particular, \(|F|\leq T\), and \(F\) has bidegree \((2,2)\).

\begin{proposition}[The pairing closure]\label{prop:Hopf-closure}
The algebra \(\cA=\R[u,v,F]=\R[R,V,F]\) is a graded Laplacian
algebra of spherical quotient dimension two. Its closure table, in
the ordered generators \((u,v,F)\), is
\begin{equation}\label{eq:Hopf-closure}
 (\Delta u,\Delta v,\Delta F)=(32,32,0),\qquad
 \bigl(\Gamma(I_i,I_j)\bigr)=
 \begin{pmatrix}
  4u&0&4F\\
  0&4v&4F\\
  4F&4F&4uv(u+v)
 \end{pmatrix}.
\end{equation}
The quartic zero set \(C=\{F=0\}\) is minimal on its regular part,
and its geometric singular set is
\begin{equation}\label{eq:Hopf-singular-set}
 Z=(\R^{16}\times\{0\})\cup(\{0\}\times\R^{16}).
\end{equation}
\end{proposition}

\begin{proof}
Equation \eqref{eq:Hopf-normalization} gives
\[
 |\nabla_xF|^2=4uv^2,\qquad
 |\nabla_yF|^2=4u^2v,\qquad \Delta F=0.
\]
Euler's identity in each block gives the remaining entries of
\eqref{eq:Hopf-closure}, and \cref{prop:MR-generator-test}
proves Laplacian closure. Surjectivity of the Hopf map gives
the full image
\[
 \{(u,v,f):u,v\geq0,\ |f|\leq uv\}.
\]
Its interior proves algebraic independence and spherical quotient
dimension two. Since $uv$ varies at fixed $R$ and $F=0$,
$\Gamma(F,F)=4Ruv\notin\R[R,F]$.

Since \(\Gamma(F,R)=8F\) and \(\Gamma(F,T)=4RF\), we obtain
\[
 \cL(F)=-\tfrac12\Gamma(F,4RT)=-8(R^2+2T)F.
\]
This proves regular minimality and identifies the critical set as
\(Z\). At \((0,y_0)\), \(y_0\ne0\), the transverse quadratic term
is \(\langle A(h(y_0))x,x\rangle\), where
\(A(a)=\sum_i a_iA_i\). Its eigenvalues are
\(\pm|y_0|^2\), each with multiplicity eight. Hence the tangent
cone is the product of a quadratic cone of signature \((8,8)\)
with \(\R^{16}\), and the point is geometrically singular.
The other axis is identical; the vertex is singular as well.
\end{proof}

\begin{remark}[A second realization of the pairing table]
The Clifford construction in \cref{sec:Clifford-composition-background}
also realizes this table. Put
\(Q=u^2+v^2-2F\). For a nine-matrix Clifford system
\(P_0,\ldots,P_8\) on \(\R^{32}\), the generators
\[
 \widehat R=|z|^2,\qquad
 \widehat V=\langle P_0z,z\rangle,\qquad
 \widehat Q=\sum_{i=0}^8\langle P_iz,z\rangle^2
\]
have the same gradient and Laplacian identities as \((R,V,Q)\).
Both images are
\(\{(R,V,Q):R\geq0,\ V^2\leq Q\leq R^2\}\).
The inverse relations
\(u=(R+V)/2\), \(v=(R-V)/2\), and \(F=(R^2+V^2-2Q)/4\)
therefore identify their graded differential algebras.
\end{remark}

\begin{theorem}[A polynomial subcalibration of the pairing cone]
\label{thm:Hopf-pairing}
The potential \(G=TF=|x|^2|y|^2F\) has degree eight and
subcalibrates \(C=\{F=0\}\). Its oriented boundary current
\(\partial\llbracket\{F<0\}\rrbracket\) is strictly area minimizing
in \(\R^{32}\). Moreover,
\begin{equation}\label{eq:Hopf-Hardy-bound}
 Q_C(\varphi)\geq\frac{449}{4}\int_C r^{-2}\varphi^2,
 \qquad \varphi\in C_c^\infty(C_{\reg}),\quad r=\sqrt R.
\end{equation}
\end{theorem}

\begin{proof}
The three blocks
\[
 E_0=(u-v)^2,\qquad E_1=2(T+F),\qquad E_2=2(T-F)
\]
have weights $(1,8,8)$. The centered generator $F/2$ gives
$h=16$, $r_0=1$, and $c=1/2$ in
\cref{thm:quartic-table-existence}. Its balanced specialization
\cref{prop:balanced-quartic-multiplier} gives $G=TF$ and
\begin{align}
 \Delta G&=40RF,\qquad
 |\nabla G|^2=4RT(T^2+3F^2),\label{eq:Hopf-gradient}\\
 \cL(G)&=32TF\bigl[T^2(2R^2-T)
                    +3F^2(4R^2-T)\bigr].\label{eq:Hopf-defect}
\end{align}
The Hardy constant in \eqref{eq:bihomogeneous-Hardy} is $449/4$
at $(N,h)=(32,16)$.
\end{proof}

\begin{remark}[Relation with the OT--FKM and Riedler constructions]
\label{rem:Hopf-FKM-Riedler}
The matrices \(B_i=\operatorname{diag}(A_i,-A_i)\) form a Clifford
system on \(\R^{32}\). Its OT--FKM polynomial
\cite[Section~1.3]{Radeschi14} is
\[
 \Psi=R^2-2\sum_{i=1}^9\langle B_i(x,y),(x,y)\rangle^2
     =4F-V^2.
\]
It has multiplicities \((8,7)\). Applying the centering formula
\eqref{eq:iso-F-z} to \(-\Psi\), with ordered multiplicities
\((7,8)\), shows that its minimal regular hypercone is
\(\{\Psi+R^2/15=0\}\). Nonvertex singularities distinguish $C$ from every cone over a
regular isoparametric hypersurface. They also distinguish it from
Riedler's family: its singular dimension is $N-2k$, so congruence
would require $N=32$, $2k=16$, but sixteen Clifford matrices
require ambient dimension at least $256$
\cite[Remark~5.2.6(1) and Table~4.1]{RiedlerThesis25}. Although
\(F=\langle\cdot,\cdot\rangle\circ(h,h)\), the product map
\(\Phi(x,y)=(h(x),h(y))\) satisfies
\[
 D\Phi D\Phi^{\mathsf T}=4\operatorname{diag}(u\Id_9,v\Id_9).
\]
It is not horizontally conformal on \(u\ne v\), so the pullback
criterion in \cref{prop:pullback-subcalibration} does not apply
to this map on \(\R^{32}\).
\end{remark}

\paragraph{The balanced section inside a Simons cone.}
\label{sec:Hopf-intrinsic-extension}

For the pairing polynomial above, the ambient space for comparison
is now itself a minimal cone. Let
\[
 \mathcal B=\{V=0\}=\{|x|=|y|\},\qquad
 D=\{V=F=0\},\qquad E_{\mathcal B}=\mathcal B\cap\{F<0\}.
\]
The Simons cone \(\mathcal B\) is a singular space of dimension
\(31\), smooth away from the origin. We use its induced metric
and the orientation determined by \(\nabla V\).
Intrinsic perimeter is variation in the smooth part, extended across
the vertex by the cutoff below. Write it as \(\Per_{\mathcal B}\), and set
\(D_a=\partial\llbracket E_{\mathcal B}\rrbracket
\mathbin{\llcorner}B_a(0)\).

\begin{theorem}[The pairing section of a Simons cone]
\label{thm:Hopf-Simons-section}
The cone \(D\) has dimension \(30\), is smooth away from the
origin, and is minimal in \(\R^{32}\). The intrinsic potential
\(G_{\mathcal B}=RF\) subcalibrates it in \(\mathcal B\).
For every \(a>0\), the integral current \(D_a\) minimizes mass
among integral \(30\)-currents supported in
\(\mathcal B\cap\overline B_a(0)\) with boundary \(\partial D_a\).
Thus \(D_a\) solves the Plateau problem in the singular ambient
space \(\mathcal B\), with its Euclidean induced area.
\end{theorem}

\begin{proof}
On \(\mathcal B\setminus\{0\}\), the gradients of \(V\) and
\(F\) are nonzero and orthogonal by \eqref{eq:Hopf-closure}.
Hence \(D\setminus\{0\}\) is smooth of codimension two.
The unit normal to \(\mathcal B\) is
\(n_{\mathcal B}=(x,-y)/\sqrt R\). The bidegree of \(F\) gives
\(\Hess F((x,-y),(x,-y))=-4F\). Since \(\mathcal B\) is minimal,
the induced operators satisfy
\begin{equation}\label{eq:Hopf-intrinsic-closure}
 \begin{gathered}
 \Delta_{\mathcal B}R=62,\qquad
 \Delta_{\mathcal B}F=4F/R,\\
 \Gamma_{\mathcal B}(R,R)=4R,\qquad
 \Gamma_{\mathcal B}(R,F)=8F,\qquad
 \Gamma_{\mathcal B}(F,F)=R^3.
 \end{gathered}
\end{equation}
Here \(\nabla F\) is tangent to \(\mathcal B\), and \(T=R^2/4\).
The same chain rule as in \cref{thm:finite-reduction} now yields
\begin{equation}\label{eq:Hopf-intrinsic-defect}
 \begin{gathered}
 \Delta_{\mathcal B}(RF)=82F,\qquad
 |\nabla_{\mathcal B}(RF)|^2=R^5+20RF^2,\\
 \cL_{\mathcal B}(RF)=16RF(2R^4+85F^2),
 \end{gathered}
\end{equation}
where \(\cL_{\mathcal B}\) is \eqref{eq:L-definition} for the
induced metric. The normalized gradient
\[
 X=\frac{\nabla_{\mathcal B}(RF)}{|\nabla_{\mathcal B}(RF)|}
\]
is smooth on \(\mathcal B\setminus\{0\}\), equals the exterior
normal along \(D\setminus\{0\}\), and its divergence has the
strict sign of \(F\) off \(D\).

We apply the proof of \cref{thm:removable-subcalibration} in the
smooth part of \(\mathcal B\). Only the vertex requires a cutoff.
For a radial cutoff \(\eta_\eps\), zero on \(B_\eps(0)\) and
one outside \(B_{2\eps}(0)\),
\[
 \int_{\mathcal B\cap B_{2\eps}(0)}
       |\nabla_{\mathcal B}\eta_\eps|\,d\cH^{31}
 =O(\eps^{30}).
\]
Hence extending a bounded characteristic function across the vertex
adds no perimeter mass. The estimate $|\diver_{\mathcal B}X|\le c/r$
from homogeneity makes the divergence error vanish as well. For every relatively compact open set \(U\subset\mathcal B\)
and every locally finite-perimeter set \(H\subset\mathcal B\) with
\(H\triangle E_{\mathcal B}\Subset U\),
\begin{equation}\label{eq:Hopf-intrinsic-comparison}
 \Per_{\mathcal B}(H;U)-\Per_{\mathcal B}(E_{\mathcal B};U)
 \geq\int_{H\triangle E_{\mathcal B}}
       |\diver_{\mathcal B}X|\,d\cH^{31}\geq0.
\end{equation}
The set \(E_{\mathcal B}\) has locally finite perimeter because
its boundary is a cone with a smooth compact link.

For the asserted integral-current comparison, cone a competitor
minus \(D_a\) to the origin inside \(\mathcal B\cap\overline B_a(0)\).
This gives an integral \(31\)-dimensional filling. On the oriented smooth part of \(\mathcal B\), write this filling as
\(\llbracket k\rrbracket_{\mathcal B}\), with integer-valued
locally BV multiplicity, and set
\(v=\one_{E_{\mathcal B}}+k\) and \(H=\{v\geq1\}\).
The integer-level decomposition used in
\cref{thm:strict-minimality-flux} gives
\(\lvert D\one_H\rvert\leq|Dv|\) on the smooth part.
Writing \(T\) for the competitor, for \(b>a\) we obtain
\[
 0\leq \Per_{\mathcal B}(H;B_b)-
          \Per_{\mathcal B}(E_{\mathcal B};B_b)
 \leq \mathbf M(T)-\mathbf M(D_a).
\]
The inequalities use \eqref{eq:Hopf-intrinsic-comparison} and
cancellation of the common exterior mass; $D$ has no mass on a sphere.
The representation and integer-level decomposition hold in smooth
charts. The vertex cutoff extends the characteristic function,
and no positive-dimensional integral current is supported at the
vertex. Thus the comparison includes competitors meeting it.

Finally, the divergence in \eqref{eq:Hopf-intrinsic-defect} vanishes
on \(D\), so \(D\) is minimal in \(\mathcal B\).
For its unit normal \(\nu\) within \(\mathcal B\), the two block
lengths satisfy \(|\nu_x|=|\nu_y|\), by the partial gradient
identities above and \(u=v\). Hence
\(\mathrm{II}_{\mathcal B}(\nu,\nu)=0\). Since
\(\operatorname{tr}_{T\mathcal B}\mathrm{II}_{\mathcal B}=0\),
its trace on \(TD\) is also zero. The second fundamental form
composition formula proves Euclidean minimality of \(D\setminus\{0\}\).
The smooth compact link gives
\(\int_{D\cap B_{2\eps}}|\nabla_D\eta_\eps|\,d\cH^{30}
=O(\eps^{29})\to0\), so the Euclidean first variation also vanishes
across the vertex.
\end{proof}

The Euclidean area-minimizing cones of codimension two in
\cite[Theorem~B]{CuiFKM25} are also area-minimizing hypersurfaces within
the Simons cones containing them. Indeed, let \(K\subset\mathcal S\)
be such a cone and put \(K_a=K\mathbin{\llcorner}B_a(0)\).
Every integral current \(T\) supported in
\(\mathcal S\cap\overline B_a(0)\) with
\(\partial T=\partial K_a\) is also an admissible Euclidean
competitor, so \(\mathbf M(K_a)\leq\mathbf M(T)\).
The induced area on \(\mathcal S\) is the Euclidean area.


